% !TEX TS-program = lualatex
\documentclass[12pt]{article}
\usepackage{fontspec}\setmainfont{MFB Oldstyle}
\usepackage{url}
\usepackage{realscripts}
\usepackage[main=english]{babel}
\usepackage{hologo}
\usepackage[autostyle=false,babel=true,english=american]{csquotes}
\usepackage[expansion=all,protrusion,babel]{microtype}
\usepackage{titlesec}\titleformat*{\section}{\bfseries}\titleformat*{\subsection}{\bfseries}
\title{The \textit{MFB Oldstyle} Font Family\\\large Version 1.0}
\date{August 5, 2024}
\author{Daniel Benjamin Miller\thanks{dbmiller@dbmiller.org}}

\begin{document}
\maketitle

\section{Description}

\textit{MFB Oldstyle} is a serif font family designed for body text. This typeface design was originally created by Morris Fuller Benton and released by American Type Founders in 1909 as \textit{Century Oldstyle}.

\section{Features}

The font family currently includes regular, \textit{italic} and \textbf{bold} fonts; no bold italic is included (though one may be added in the future). All the fonts include superior and inferior figures, accessible both through OpenType features and the provided \texttt{mfb-oldstyle} \hologo{LaTeX} package. \textsc{Small capitals} and old-style figures ({\addfontfeatures{Numbers=OldStyle}0123456789}) are currently provided for the regular font only. For a view of character coverage, see the included character map PDF.

\section{Usage with OpenType-capable \hologo{TeX} engines}

I recommend using the OpenType versions of \textit{MFB Oldstyle}, along with a modern \hologo{TeX} engine (such as \hologo{LuaTeX} or \hologo{XeTeX}). If you are using \hologo{LuaLaTeX} or \hologo{XeLaTeX}, use \texttt{fontspec} to select the font and the appropriate options, e.g., by adding \texttt{\textbackslash usepackage\{fontspec\}\textbackslash setmainfont\{MFB Oldstyle\}} to your preamble. For detailed options, see that package's documentation.

\section{Usage with legacy \hologo{TeX} engines}

If you wish to use the fonts with a legacy \hologo{TeX} engine, Type 1 fonts and the appropriate support files are provided. To set use the Type 1 version of \textit{MFB Oldstyle} in \hologo{LaTeX}, use the following command:

\begin{verbatim}
\usepackage{mfb-oldstyle}
\end{verbatim}

By default, tabular lining figures (0123456789) are used. The package-level options \texttt{tabular} and \texttt{proportional} can be used to switch between tabular ({\addfontfeatures{Numbers=Tabular}0123456789}) and proportional ({\addfontfeatures{Numbers=Proportional}0123456789}) versions of lining figures. Old-style figures ({\addfontfeatures{Numbers=OldStyle}0123456789}) can technically be set as the default using package option \texttt{oldstyle}, but, since old-style figures are only provided for the regular font, this will disable use of the italic and bold fonts by default. (For a more graceful fallback, try using a modern \hologo{TeX} engine). When using the legacy \hologo{LaTeX} package, superior figures can be accessed on a one-off basis using the text command \texttt{textsu}, and inferior figures can be accessed using the text command \texttt{textin}.

\section{Copying}

The \textit{MFB Oldstyle} fonts have been placed in the public domain and can be used without restriction. For a detailed copyright waiver with explicit fallback provisions, see \texttt{COPYING}.

\section{Source files and development}

Source files for the \textit{MFB Oldstyle} fonts are available on GitHub.\footnote{\url{https://github.com/dbenjaminmiller/mfb-oldstyle}} Suggestions for feature additions are welcome on GitHub (preferred) or via email.



\end{document}